% ============================================================
%  N.1.1  Architectures
% ============================================================
\subsection{Architectures}
\label{sec:architectures}

The two spatial architectures selected for the experiments are
DepFiN~\cite{depfin_paper} and Eyeriss~v1~\cite{eyeriss_paper}.
Table~\ref{tab:orig-arch-comparison} summarises the key parameters of the
\emph{original} designs as published in the respective papers.

% ------------------------------------------------------------------
%  Comparison table: original DepFiN vs original Eyeriss v1
% ------------------------------------------------------------------
\begin{table}[htbp]
  \centering
  \caption{Comparison of the original DepFiN and Eyeriss~v1 architectures.}
  \label{tab:orig-arch-comparison}
  \small
  \begin{tabular}{l l l}
    \toprule
    \textbf{Parameter}          & \textbf{DepFiN}                      & \textbf{Eyeriss v1}                 \\
    \midrule
    Technology node              & 22\,nm                               & 65\,nm                              \\
    Clock frequency              & $\sim$930\,MHz                       & 200\,MHz                            \\
    Dataflow                     & Depth-first (weight-stationary)      & Row-stationary                      \\
    PE array                     & $16 \times 128 = 2048$               & $12 \times 14 = 168$                \\
    \addlinespace
    On-chip memory               & Split FMEM + WMEM                    & Unified GlobalBuffer                \\
    FMEM / GlobalBuffer size     & 1056\,KB                             & 128\,KB                             \\
    FMEM wordline                & 1056\,bit (2 banks)                  & 64\,bit                             \\
    WMEM size                    & 524\,KB                              & ---                                 \\
    WMEM wordline                & 128\,bit (4 banks)                   & ---                                 \\
    \addlinespace
    PE registers                 & WeightReg (384),                     & InputSpad (24),                     \\
                                 & InputReg (24),                       & WeightSpad (384),                   \\
                                 & AccumulationReg (32)                 & PsumSpad (32)                       \\
    Accumulator precision        & 32\,bit                              & 16\,bit                             \\
    Operand precision            & 8\,bit                               & 16\,bit                             \\
    \addlinespace
    DRAM type                    & LPDDR4                               & LPDDR4                              \\
    DRAM bandwidth               & 12\,B/cycle                          & 8\,B/cycle                          \\
    FMEM read bandwidth          & 132\,B/cycle                         & ---                                 \\
    GB bandwidth                 & ---                                  & 32\,B/cycle                         \\
    \addlinespace
    Tile size                    & Fixed $1 \times 128$ patch            & Determined by mapping               \\
    Spatial mapping (cols)       & Output width ($Q$)                   & Output width ($Q$) + channels ($M$) \\
    Spatial mapping (rows)       & Output channels ($M$)                & Filter ($S$) + channels ($C$, $M$)  \\
    \bottomrule
  \end{tabular}
\end{table}

Both architectures were modelled in the FactorFlow framework based on the
parameters published in~\cite{depfin_paper,eyeriss_paper}.
However, because the systolic arrays must support a variety of fused CNN
workloads---not only the single fixed-size layer for which each original design
was optimised---several modifications were necessary.
We therefore refer to the modified variants as \textbf{DepFiN-Like} and
\textbf{Eyeriss-Like} throughout this chapter.

% ------------------------------------------------------------------
\subsubsection{DepFiN-Like}
\label{sec:depfin-like}

DepFiN is a depth-first CNN accelerator with a split on-chip memory
hierarchy: a \emph{Feature Memory} (FMEM) that stores input and output
activations while bypassing weights, and a \emph{Weight Memory} (WMEM) that
stores weights while bypassing all activation operands.
In multi-layer (fused) mode, intermediate activations between consecutive
layers are kept entirely on-chip---DRAM bypasses the intermediate operands
(\texttt{int\_in}, \texttt{int\_out})---so that inter-layer data never
touches external memory.
Each fused layer receives its own pair of spatial fanout levels:
\texttt{SACols} distributes the output spatial width ($Q$ or $X_i$) across
the PE columns, while \texttt{SARows} distributes the output channels
($M$ or $Z_i$) across the PE rows.

Compared with the original DepFiN design (see
Section~\ref{sec:validation}), the following modifications were applied:

\begin{enumerate}
  \item \textbf{Configurable tile size.}
        The original DepFiN fixes the output tile to a $1 \times 128$
        pixel patch that fills the 128 PE columns exactly.
        Our version relaxes this constraint, allowing the output tile
        size to be set independently for each workload.
        This is essential for exploring the effect of tile size on
        energy, latency, and DRAM traffic across workloads with
        different spatial dimensions (see Section~\ref{sec:cs1-tile}).

  \item \textbf{Variable PE array dimensions.}
        Instead of the fixed $16 \times 128$ array, we sweep the number
        of PE rows and columns independently
        (Sections~\ref{sec:cs2-rows}--\ref{sec:cs3-cols})
        and explore different aspect ratios at a constant total PE count
        (Section~\ref{sec:cs4-aspect}).

  \item \textbf{Feature Memory bandwidth.}
        The FMEM read bandwidth is a fixed physical interface parameter
        (132\,B/cycle, determined by the 1056-bit wordline).
        It does \emph{not} scale automatically with the PE array size or
        the tile size.
        An optional bandwidth-scaling mode is available---in which the
        bandwidth is scaled proportionally to
        $\mathit{tile\_size} / 128$---and is used in the tile-size
        sensitivity analysis (Section~\ref{sec:cs1-tile}).

  \item \textbf{Mapping.}
        The depth-first dataflow is preserved: weights are loaded into
        WMEM for all fused layers, and activations stream through FMEM
        tile-by-tile.
        The mapping is determined by assigning spatial factors greedily:
        \texttt{SACols} receives the highest divisor of the output tile
        width that fits in the column mesh, and \texttt{SARows} receives
        the highest divisor of the output channel count that fits in the
        row mesh.
        Remaining weight-dimension factors ($C$, $R$, $S$) are assigned
        to WMEM, and residual tiling factors are pushed to DRAM.
        When the tile size or PE array dimensions change across case
        studies, the mapping adapts accordingly.
\end{enumerate}

% ------------------------------------------------------------------
\subsubsection{Eyeriss-Like}
\label{sec:eyeriss-like}

Eyeriss~v1 is a row-stationary accelerator with a unified
\emph{GlobalBuffer} (128\,KB) that stores input activations and output
partial sums while bypassing weights (weights are delivered directly
from DRAM to the per-PE weight registers).
In the original design each PE contains three register files:
an \emph{InputSpad} for input activations,
a \emph{WeightSpad} for filter weights,
and a \emph{PsumSpad} for partial-sum accumulation.
The row-stationary dataflow is expressed through the spatial fanout:
\texttt{SACols} distributes output spatial positions ($Q$) and output
channels ($M$) across the 14 PE columns, while \texttt{SARows}
distributes filter rows ($S$), input channels ($C$), and output channels
($M$) across the 12 PE rows.
Each PE therefore processes a 1-D convolution row: filter weights remain
stationary in the weight register while input activations stream through,
and partial sums accumulate across input channels along the columns.

Compared with the original Eyeriss, the following modifications were
applied:

\begin{enumerate}
  \item \textbf{Intermediate register for fusion.}
        The original Eyeriss has no mechanism for passing intermediate
        activations between fused layers on-chip.
        We add a fourth per-PE register file, the
        \emph{IntermediateOutputRegister}, which stores intermediate
        activations produced by one layer and consumed by the next.
        This register bypasses all other operands (inputs, weights,
        final outputs) and, together with the DRAM bypass of
        \texttt{int\_in}/\texttt{int\_out}, enables depth-first
        layer fusion on Eyeriss
        (see Section~\ref{sec:partial-full-fusion}).

  \item \textbf{Variable PE array dimensions.}
        With only 168~PEs ($14 \times 12$), each PE must process a
        disproportionately large share of the workload.
        For weight-dominant networks such as VGG16 and ResNet18, this
        forces the per-PE weight register to hold thousands of entries,
        which is physically impractical.
        We therefore explore four total PE counts:
        2\,048, 4\,096, 8\,192, and 16\,384.
        Going above 16\,384 would exceed realistic budgets for
        low-power edge accelerators; going below 4\,096 pushes the
        required weight-register size to several thousand entries per PE,
        making the design infeasible.
        The sensitivity of register sizes and EDP to the PE count and
        aspect ratio is analysed in
        Sections~\ref{sec:cs1-wreg}--\ref{sec:cs5-aspect}.

  \item \textbf{Register sizing.}
        Because the Eyeriss feasibility problem is multi-dimensional
        (InReg, WReg, IntermediateReg, and OutReg must all be
        satisfiable simultaneously for each workload and PE
        configuration), the per-PE register sizes are treated as
        experimental variables.
        Case studies CS1--CS3 sweep WReg, IntReg, and OutReg in turn
        while searching for the minimum feasible PE configuration at
        each size (Section~\ref{sec:cs1-wreg}--\ref{sec:cs2cs3-regs}).

  \item \textbf{Mapping.}
        The row-stationary dataflow is preserved.
        Spatial factors are assigned greedily:
        \texttt{SACols} first allocates the highest divisor of the
        output tile width, then assigns remaining columns to output
        channels~$M$;
        \texttt{SARows} assigns factors to filter rows~$S$, then input
        channels~$C$, then output channels~$M$.
        The GlobalBuffer stores the residual tiling of activations
        (with $P = 1$), and DRAM handles the outermost loops.
        As with DepFiN-Like, the mapping adapts when tile sizes or PE
        dimensions change across experiments.
\end{enumerate}

Table~\ref{tab:modified-arch-comparison} summarises the key differences
between the original and modified versions of both architectures.

% ------------------------------------------------------------------
%  Comparison table: original vs modified
% ------------------------------------------------------------------
\begin{table}[htbp]
  \centering
  \caption{Summary of modifications applied to the original architectures.}
  \label{tab:modified-arch-comparison}
  \small
  \begin{tabular}{l p{5cm} p{5cm}}
    \toprule
    \textbf{Aspect}
      & \textbf{DepFiN $\to$ DepFiN-Like}
      & \textbf{Eyeriss $\to$ Eyeriss-Like} \\
    \midrule
    PE array
      & $16 \times 128$ (fixed) $\to$ variable rows/cols
      & $12 \times 14$ (fixed) $\to$ 2\,048--16\,384 total PEs \\
    \addlinespace
    Tile size
      & Fixed $1 \times 128$ $\to$ configurable per workload
      & Determined by mapping $\to$ configurable per workload \\
    \addlinespace
    Registers
      & Unchanged default sizes
      & Added IntermediateOutputReg (320 entries);
        WReg/InReg/OutReg sizes swept as variables \\
    \addlinespace
    FMEM / GB bypass
      & Weights bypassed (unchanged)
      & Weights bypassed (unchanged) \\
    \addlinespace
    DRAM bypass (fusion)
      & \texttt{int\_in}, \texttt{int\_out} bypassed
      & \texttt{int\_in}, \texttt{int\_out} bypassed \\
    \addlinespace
    Dataflow
      & Depth-first preserved
      & Row-stationary preserved \\
    \addlinespace
    Technology
      & 22\,nm, $0.5\times$ scale
      & 22\,nm, $0.5\times$ scale (orig.\ 65\,nm) \\
    \bottomrule
  \end{tabular}
\end{table}






My version:
The spatial architectures chosen for the experiments are DepFin [citation to the paper], Eyeriss [citation to the paper] ("refer to the table of comparison"), we modeled both architectures based on the characteristics and parameters published in [citations to both paper], but since the SAs must be able to perform the computation of a variety of fused CNN workloads, we have manually performed changes on both of architectures, so we will refer to them to DepFin-Like and Eyeriss-Like.


DepFin-Like: with respect to the Section ("refer to the correct section in validation"), our version relax the constraint of having a patch of 1 pixel height and 128 pixel wide that is corresponding to the tile size used, this was particularly important in order to able the exploration of different tile sizes per different workloads. We relaxed the constraint of having a structure of 16x128 PEs, we created more versions of DepFin analyzing how changing the number of PEs will or not will affects to the computation results and exploring the aspect ratio of PEs for each workload. We respected the DepthFirst ("refer to the section in background") nature of DepFin, preserving the dataflow,  but with respect to a classic DepFin we will change the mapping thanks to the different tiling adopted through the chosen of tile size.



Eyeriss-Like: with respect to the ("cite: paper") our version of Eyeriss must address not only to SL/LBL processing ("refer to the section in background"), but fusion, in order to do so we introduced a intermediate register, the unique maintaining the intermediate input and output data and bypassing all the other operands ("refer to the section in background").  We respected the row stationary nature of Eyeriss, maintaing the rows in weight register. 
Since we respected the row stationary nature of Eyeriss ("refer to the section in the background"), we had to relax the constraint of having a structure 14x12 PEs, with only 168 PEs (14×12), each PE must process a larger fraction of the workload, meaning the WRegister per PE must hold more weight values. For large workloads like VGG16/ResNet18, this makes the register infeasibly large. We created more version of Eyeriss with 4 possible total of PEs: 2048, 4096, 8192, 16184, going above 16384 is unrealistic for low-power edge accelerators; going below 4096 forces WReg to several thousand entries, which is impractical. We will analyze how changing the number of PEs will or will not affects the required sizes of registers per workload and exploring the aspect ratio of PEs for each workload. As stated, we preserved the row stationary dataflow  but we will change the mapping thanks to the different tiling adopted through the chosen of tile size.
